% ============================================================
% Chapter 6: Experimental Design
%   Section: Construction of the Custom Validation Set
%     Subsection: Tie-Break Rules
% Included from chapters/06-experimental-design/03-custom-validation-set/custom-validation-set.tex
% ============================================================

\subsection{Tie-Break Rules}
\label{subsec:custom-rules}

The decisions that cause the most hesitation during annotation are
also those on which a single annotator is least consistent with
themself over time. The following four rules were therefore fixed
before labelling began and are applied without exception; the
\texttt{annotatorNote} records which rule decided a close call.

\paragraph{Rule 1 --- The organisation's own source.}
Positive-news claims frequently originate in the press release of
the very NGO, company or institution the article is about. A claim
for which the \emph{only} available evidence is that organisation's
own communication (press release, report, website or social media
post) is labelled \emph{not enough evidence}
(\texttt{UNVERIFIED}), and the source is named in the note. An
organisation reporting on itself is not independent confirmation,
and a claim the system could only confirm by reading the claimant's
own release is not one it has verified.

\paragraph{Rule 2 --- Numeric tolerance.}
A figure is considered supported when it lies within $\pm10\%$ in
relative terms of the figure in the source (an article's ``30\%''
against a source's ``28\%'' is supported), or when it is a reasonable
rounding of it (``nearly 3,000'' for 2,870). Beyond that margin, the
claim is \emph{partially supported} if the figure is a detail of an
otherwise correct claim, and \emph{refuted} if the figure is the
claim.

\paragraph{Rule 3 --- Date of the evidence.}
Only evidence available on the article's publication date is
admissible, as in AVeriTeC \cite{schlichtkrull2023averitec}, which
restricts evidence to documents published before the claim to avoid
temporal leakage. A claim that could not be checked when it was
published and was confirmed later remains \emph{not enough evidence};
the verdict does not change. Evidence that did not yet exist cannot
have made a claim well-founded when it was made, and fact-checking
evaluations that ignore this condition have been shown to overstate
what a system could do in practice \cite{glockner2022missing}.

\paragraph{Rule 4 --- Hierarchy of sources.}
Sources are ranked as follows: primary source (a report, official
statistic or peer-reviewed study) $>$ reference media outlet $>$
press release $>$ social media. When sources conflict, the higher
tier prevails; when sources of the same tier conflict and none is
decisive, the claim is labelled \emph{conflicting evidence}
(\texttt{MISLEADING}). The tier of the strongest source used is
recorded in \texttt{sourceTier}, which also allows the evaluation to
be broken down by how well-sourced each claim was.
